\noindent
Let $R \to S$ be a ring map. Denote $R[S]$ the polynomial ring
whose variables are the elements $s \in S$. Let's denote $[s] \in R[S]$
the variable corresponding to $s \in S$. Thus $R[S]$ is a free
$R$-module on the basis elements $[s_1] \ldots [s_n]$ where
$s_1, \ldots, s_n$ ranges over all unordered sequences of elements of $S$.
There is a canonical surjection
\begin{equation}
\label{equation-canonical-presentation}
R[S] \longrightarrow S,\quad [s] \longmapsto s
\end{equation}
whose kernel we denote $I \subset R[S]$. It is a simple observation that
$I$ is generated by the elements
$[s + s'] - [s] - [s']$, $[s][s'] - [ss']$ and $[r] - r$.
According to Lemma \ref{lemma-differential-seq}
there is a canonical map
\begin{equation}
\label{equation-naive-cotangent-complex}
I/I^2 \longrightarrow \Omega_{R[S]/R} \otimes_{R[S]} S
\end{equation}
whose cokernel is canonically isomorphic to $\Omega_{S/R}$. Observe that
the $S$-module $\Omega_{R[S]/R} \otimes_{R[S]} S$ is free on the generators
$\text{d}[s]$.
